% !TeX root = ../../tesi.tex

\section{Innovazione e intelligenza artificiale}
\label{sec:innovazione}

Halue pratica uno \emph{scouting} tecnologico continuo e dedica una quota del tempo a progetti di
ricerca e sviluppo, tra cui quelli di intelligenza artificiale (\figref{fig:ricerca}). I dipendenti
dedicano allo studio una parte rilevante del proprio tempo.

\begin{figure}[htbp]
	\centering
	\setstretch{1}
	% La quota del tempo di Halue dedicata a ricerca e sviluppo, come grafico a torta.
% La fetta in bordeaux vale il 10%, il minimo dichiarato dall'azienda.

\begin{tikzpicture}[every node/.style={text=FiguraTesto}]
	\def\raggio{15mm}
	% Il resto del tempo, poi la fetta di ricerca e sviluppo (36 gradi, il 10%), staccata di poco.
	\fill[TabellaIntestazione, draw=FiguraBordo] (0, 0) -- (18:\raggio) arc[start angle=18,
		end angle=342, radius=\raggio] -- cycle;
	\fill[FiguraAccento, draw=white, line width=0.6pt, shift={(0:1mm)}] (0, 0) -- (-18:\raggio)
		arc[start angle=-18, end angle=18, radius=\raggio] -- cycle;

	% Etichette, a destra della torta.
	\draw[FiguraTenue] (0:\raggio + 1.5mm) -- (0:\raggio + 7mm);
	\node[anchor=west, align=left, font=\fontsize{7.5}{9}\selectfont]
		at (\raggio + 8mm, 0) {\voce{Ricerca e sviluppo}{oltre il 10\% del tempo}};
	\draw[FiguraTenue] (150:\raggio * 0.55) -- (150:\raggio + 7mm) -- ++(-5mm, 0);
	\node[anchor=east, align=right, font=\fontsize{7.5}{9}\selectfont]
		at ($(150:\raggio + 7mm) + (-6mm, 0)$) {\voce{Progetti per i clienti}{e altre attività}};
\end{tikzpicture}

	\caption{La quota del tempo che Halue dedica alla ricerca e sviluppo.}
	\label{fig:ricerca}
\end{figure}

L’intelligenza artificiale è il campo in cui questa propensione si vede di più. Il \emph{team}
segue da vicino le novità del settore, testa regolarmente nuovi strumenti di automazione basati
sull’IA e li porta nel lavoro di ogni giorno. Nello sviluppo l’IA esegue una prima revisione
automatica del codice e analizza i \emph{changelog} delle tecnologie in uso. Per i processi
ricorrenti l’azienda usa agenti di IA guidati da \gloss{skill}{\emph{skill}} interne, cioè
documenti di istruzioni, a volte accompagnati da \emph{script} e modelli, che descrivono come
svolgere un processo e che l’agente consulta quando il compito lo richiede. In tutti questi usi vale la stessa regola: l’IA propone
e la persona decide (\figref{fig:ciclo-ia}).

\begin{figure}[htbp]
	\centering
	\setstretch{1}
	% Il lavoro con gli strumenti di IA in Halue, come diagramma di attività UML: l'IA produce una
% proposta, la persona la rivede e, finché non è adeguata, la corregge con nuove indicazioni.
% Il disegno è largo quanto il testo (150 mm): si include con \input, senza ridimensionarlo.

\begin{tikzpicture}[
	every node/.style={text=FiguraTesto},
	azione/.style={draw=FiguraTenue, fill=white, rounded corners=3pt, text width=#1-2mm,
		minimum width=#1, minimum height=27mm, align=center, inner sep=2.5pt,
		font=\fontsize{7.5}{10}\selectfont,
		execute at begin node={\hyphenpenalty=10000\exhyphenpenalty=10000}},
	azione/.default=24mm,
	decisione/.style={draw=FiguraTenue, fill=white, diamond, minimum size=8mm, inner sep=0pt},
	iniziale/.style={circle, fill=FiguraTesto, minimum size=3mm, inner sep=0pt},
	finale/.style={circle, draw=FiguraTesto, fill=white, minimum size=3.6mm, inner sep=0pt,
		path picture={\fill[FiguraTesto] (path picture bounding box.center) circle (1.05mm);}},
	guardia/.style={font=\fontsize{7.5}{9}\selectfont, text=FiguraTenue, inner sep=1pt},
]
	\node[iniziale] (inizio) at (2mm, 0) {};
	\node[azione=26mm] (compito) at (20mm, 0)
		{\voce{Compito}{codice, documentazione, \emph{ticket}, rapporti, stime, grafica}};
	\node[azione] (proposta) at (51mm, 0)
		{\voce{Proposta dell’IA}{bozza prodotta dall’assistente}};
	\node[azione] (revisione) at (80mm, 0)
		{\voce{Revisione}{della persona: correttezza e qualità}};
	\node[decisione] (adeguata) at (102mm, 0) {};
	\node[azione] (risultato) at (126mm, 0)
		{\voce{Risultato}{pronto per il progetto}};
	\node[finale] (fine) at (146mm, 0) {};
	\node[azione=36mm, minimum height=17mm] (indicazioni) at (76.5mm, -48mm)
		{\voce{Nuove indicazioni}{correzioni e richieste più precise}};

	\draw[freccia] (inizio) -- (compito);
	\draw[freccia] (compito) -- (proposta);
	\draw[freccia] (proposta) -- (revisione);
	\draw[freccia] (revisione) -- (adeguata);
	\draw[freccia] (adeguata) -- node[guardia, above] {[sì]} (risultato);
	\draw[freccia] (risultato) -- (fine);
	\node[guardia, above=1mm of adeguata] {È adeguata?};
	\draw[freccia] (adeguata.south) |- node[guardia, pos=0.25, right] {[no]} (indicazioni.east);
	\draw[freccia] (indicazioni.west) -| (proposta.south);
\end{tikzpicture}

	\caption{Il ciclo di lavoro con gli strumenti di intelligenza artificiale in Halue.}
	\label{fig:ciclo-ia}
\end{figure}

L’azienda riscontra una riduzione dei tempi, ma non dispone ancora di misure precise su questo
beneficio e sull’impatto sulla qualità. Nel \emph{team} è condivisa un’opinione su cui torno nella
\secref{sec:maturazione}: nelle mani di un programmatore esperto l’IA è un buon programmatore
\emph{junior}, ma non può ancora sostituire un \emph{senior}.

Nei progetti di IA per i clienti Halue adotta alcune regole che valgono anche per il mio progetto:

\begin{itemize}
	\item i dati vengono sempre dal catalogo reale, tramite strumenti e ricerca, senza affidarsi alla
	memoria del modello;
	\item le azioni con effetti esterni restano sotto controllo umano;
	\item la qualità delle risposte si valuta con \emph{test} manuali su casi d’uso reali;
	\item il modello si sceglie in base al compito, in modo da tenere sotto controllo il consumo dei
	\gloss{token}{\emph{token}}, i frammenti di testo in cui il modello scompone ciò che legge e
	scrive e su cui si calcola il costo del suo utilizzo.
\end{itemize}

Questa propensione all’innovazione passa anche dagli \emph{stage}, che Halue usa per esplorare temi
che i progetti per i clienti non lasciano il tempo di approfondire. Approfondisco questo ruolo nella
\secref{sec:stage-strategia}.
